\lecture{5}{Rotman §2.6}{Quotient Groups}

\section{Quotient Groups}

The group \(\mathbb Z_m\) of residue classes was built from \(\mathbb Z\) and its subgroup
\(m\mathbb Z\). This construction is the prototype of a general way of building new groups
from given ones, called \emph{quotient groups}. Along the way, group theory gives new proofs
of the theorems of Fermat, Euler, and Wilson.

\subsection{The Integers mod \(m\), Revisited}

Fix \(m \geq 2\). We constructed the residue classes \([a]_m = a + m\mathbb Z\) and the set
\(\mathbb Z_m\) in Lecture~2. Rotman's Propositions 2.98--2.103 are results we have
already proved:
\begin{itemize}
	\item \([a]_m = [b]_m\) if and only if \(a \equiv b \pmod m\)
	      (\autoref{prop:residue-class-equality}); in particular \([a]_m = [0]_m\) if and only
	      if \(m \mid a\);
	\item \(\mathbb Z_m = \{[0]_m, [1]_m, \dots, [m-1]_m\}\) has exactly \(m\) elements
	      (\autoref{cor:Zn-list});
	\item \([a]_m + [b]_m = [a+b]_m\) is well defined (\autoref{lem:addition-well-defined}),
	      and \(\mathbb Z_m\) is an abelian group of order \(m\) under it
	      (\autoref{thm:Zn-additive-group});
	\item \([a]_m[b]_m = [ab]_m\) is well defined (\autoref{lem:multiplication-well-defined});
	      it is associative and commutative with identity \([1]_m\), because these laws are
	      inherited from \(\mathbb Z\) exactly as for addition.
\end{itemize}
The group axioms in \(\mathbb Z_m\) are ``inherited'' from those of \(\mathbb Z\). When the
modulus is clear from context, we sometimes drop the subscript and write \([a]\) for
\([a]_m\).

\begin{proposition}\label{prop:Zm-cyclic}
	\(\mathbb Z_m\) is a cyclic group of order \(m\) with generator \([1]_m\). Consequently,
	every cyclic group of order \(m\) is isomorphic to \(\mathbb Z_m\).
\end{proposition}

\begin{proof}
	If \(0 \leq r < m\), then \([r]_m = [1]_m + \cdots + [1]_m\) (\(r\) summands), so
	\(\mathbb Z_m = \langle[1]_m\rangle\). The second statement follows from
	\autoref{eg:cyclic-iso}.
\end{proof}

\begin{remark}
	Here is an alternative construction. Let \(G_m = \{0, 1, \dots, m-1\}\), with
	\[
		a \oplus b = \begin{cases}
			a + b     & \text{if } a + b \leq m-1, \\
			a + b - m & \text{if } a + b > m-1.
		\end{cases}
	\]
	Although this definition looks simpler, proving associativity is now very tedious, and
	proofs with it usually require case analyses (compare \autoref{eg:cyclic-iso}).
\end{remark}

\(\mathbb Z_m\) is \emph{not} a group under multiplication, since some elements, e.g.\
\([0]_m\), have no inverse. We shall repeatedly need two consequences of B\'ezout's identity.

\begin{lemma}\label{lem:coprime-divisibility}
	Let \(a, b, m, n\) be integers.
	\begin{enumerate}[(i)]
		\item \emph{(Euclid's lemma)} If \(\gcd(m, a) = 1\) and \(m \mid ab\), then
		      \(m \mid b\). In particular, if \(p\) is prime and \(p \mid ab\), then
		      \(p \mid a\) or \(p \mid b\).
		\item If \(\gcd(m, n) = 1\), \(m \mid a\), and \(n \mid a\), then \(mn \mid a\).
	\end{enumerate}
\end{lemma}

\begin{proof}
	(i) By \autoref{cor:bezout-criterion} there are \(x, y\) with \(mx + ay = 1\). Then
	\(b = mbx + (ab)y\), and \(m\) divides both terms. If \(p\) is prime and \(p \nmid a\),
	then \(\gcd(p, a) = 1\), since the only positive divisors of \(p\) are \(1\) and \(p\).

	(ii) Write \(a = mq\). Then \(n \mid mq\) with \(\gcd(n, m) = 1\), so \(n \mid q\) by
	(i), and \(mn \mid mq = a\).
\end{proof}

\begin{proposition}\label{prop:Zp-units}
	\begin{enumerate}[(i)]
		\item If \(\gcd(a, m) = 1\), then \([a]_m[x]_m = [b]_m\) can be solved for
		      \([x]_m \in \mathbb Z_m\).
		\item If \(p\) is prime, then \(\mathbb Z_p^\times\), the set of nonzero elements of
		      \(\mathbb Z_p\), is a multiplicative abelian group of order \(p - 1\).
	\end{enumerate}
\end{proposition}

\begin{proof}
	(i) Choose \(s, t\) with \(sa + tm = 1\). Then \([s]_m[a]_m = [1]_m\), so
	\([x]_m = [sb]_m\) is a solution: \([a]_m[sb]_m = [s]_m[a]_m[b]_m = [b]_m\).

	(ii) If \(0 < a < p\), then \(\gcd(a, p) = 1\). Hence \(\mathbb Z_p^\times\) coincides
	with the group \(\mathbb Z_p^*\) of units of \autoref{thm:units-group}, and it has
	\(p - 1\) elements, being \(\mathbb Z_p\) with \([0]_p\) removed.
\end{proof}

We will prove later (Rotman, Theorem 3.122) that \(\mathbb Z_p^\times\) is a \emph{cyclic}
group of order \(p - 1\). Here is a new proof of Fermat's theorem, entirely different from
the one in \autoref{eg:mckay}.

\begin{corollary}[Fermat]\label{cor:fermat}
	If \(p\) is a prime and \(a \in \mathbb Z\), then \(a^p \equiv a \pmod p\).
\end{corollary}

\begin{proof}
	It suffices to show that \([a^p]_p = [a]_p\). If \([a]_p = [0]_p\), then
	\([a^p]_p = [a]_p^p = [0]_p = [a]_p\). If \([a]_p \neq [0]_p\), then
	\([a]_p \in \mathbb Z_p^\times\), a group of order \(p - 1\), so
	\([a]_p^{p-1} = [1]_p\) by \autoref{cor:am-equals-1}. Multiplying by \([a]_p\) gives
	\([a^p]_p = [a]_p^p = [a]_p\).
\end{proof}

If \(m\) is not prime, then \(\mathbb Z_m^\times = \mathbb Z_m - \{[0]_m\}\) is \emph{not} a
group: if \(m = ab\) with \(1 < a, b < m\), then \([a]_m, [b]_m \neq [0]_m\), but
\([a]_m[b]_m = [m]_m = [0]_m\). The right analogue of \(\mathbb Z_p^\times\) is the group
\(\mathbb Z_m^*\) of units from Lecture~2, which Rotman denotes by \(U(\mathbb Z_m)\).

\begin{definition}[Euler \(\phi\)-function]
	For \(m \geq 1\), \(\phi(m)\) is the number of integers \(k\) with \(1 \leq k \leq m\) and
	\(\gcd(k, m) = 1\).
\end{definition}

\begin{lemma}\label{lem:units-order}
	\(\mathbb Z_m^*\) is exactly the set of classes in \(\mathbb Z_m\) having a
	multiplicative inverse, and it is a multiplicative abelian group of order \(\phi(m)\).
\end{lemma}

\begin{proof}
	The first statement is \autoref{cor:unit-criterion}, and \(\mathbb Z_m^*\) is an abelian
	group by \autoref{thm:units-group}. By \autoref{cor:Zn-list}, \(\mathbb Z_m^*\) consists
	of the \([r]_m\) with \(0 \leq r \leq m-1\) and \(\gcd(r, m) = 1\), and these are
	distinct. Since \(\gcd(0, m) = m = \gcd(m, m) \neq 1\), replacing the range
	\(0 \leq r \leq m-1\) by \(1 \leq r \leq m\) does not change the count, which is
	therefore \(\phi(m)\).
\end{proof}

If \(p\) is prime, then \(\phi(p) = p - 1\) and \(\mathbb Z_p^* = \mathbb Z_p^\times\).

\begin{theorem}[Euler]\label{thm:euler}
	If \(\gcd(r, m) = 1\), then \(r^{\phi(m)} \equiv 1 \pmod m\).
\end{theorem}

\begin{proof}
	If \(G\) is a finite group of order \(n\), then \(x^n = 1\) for all \(x \in G\)
	(\autoref{cor:am-equals-1}). Here \([r]_m \in \mathbb Z_m^*\), a group of order
	\(\phi(m)\) by \autoref{lem:units-order}, so \([r]_m^{\phi(m)} = [1]_m\). In congruence
	notation, \(r^{\phi(m)} \equiv 1 \pmod m\).
\end{proof}

\begin{eg}[\(\mathbb Z_8^*\) and \(\mathbb Z_{10}^*\)]\label{eg:units-8-10}
	We have
	\[
		\mathbb Z_8^* = \{[1], [3], [5], [7]\} \cong \mathbf V,
	\]
	for \([3]^2 = [9] = [1]\), \([5]^2 = [25] = [1]\), and \([7]^2 = [49] = [1]\), so every
	nonidentity element has order \(2\) (see \autoref{prop:order-4}). Moreover,
	\[
		\mathbb Z_{10}^* = \{[1], [3], [7], [9]\} \cong \mathbb Z_4,
	\]
	for \([3]^4 = [81] = [1]\) while \([3]^2 = [9] = [-1] \neq [1]\); so \([3]\) has order
	\(4\), and \(\mathbb Z_{10}^*\) is cyclic.
\end{eg}

\begin{theorem}[Wilson]\label{thm:wilson}
	An integer \(p \geq 2\) is a prime if and only if
	\[
		(p-1)! \equiv -1 \pmod p.
	\]
\end{theorem}

\begin{proof}
	Assume \(p\) is prime. If \(a_1, \dots, a_n\) is a list of all the elements of a finite
	abelian group \(G\), then the product \(a_1a_2\cdots a_n\) equals the product of all
	elements \(a\) with \(a^2 = 1\), for every other element cancels against its inverse
	(\(a \neq a^{-1}\) pairs off with \(a^{-1}\)). In \(\mathbb Z_p^\times\), if
	\([a]^2 = [1]\), then \(p \mid a^2 - 1 = (a-1)(a+1)\), so \(p \mid a - 1\) or
	\(p \mid a + 1\) by \autoref{lem:coprime-divisibility}(i); that is, \([a] = [\pm 1]\).
	Hence the product of all elements of \(\mathbb Z_p^\times\), namely
	\([1][2]\cdots[p-1] = [(p-1)!]\), is \([1][-1] = [-1]\) if \(p\) is odd, and
	\([1] = [-1]\) if \(p = 2\). Therefore \((p-1)! \equiv -1 \pmod p\).

	Conversely, suppose \((m-1)! \equiv -1 \pmod m\), say \((m-1)! + 1 = km\). Then every
	common divisor of \(m\) and \((m-1)!\) divides \(1\), so \(\gcd(m, (m-1)!) = 1\). If
	\(m\) were composite, there would be \(a \mid m\) with \(1 < a \leq m-1\); then
	\(a \mid (m-1)!\), and \(a > 1\) would be a common divisor of \(m\) and \((m-1)!\), a
	contradiction. Therefore \(m\) is prime.
\end{proof}

\begin{remark}
	Wilson's theorem generalizes the way Euler's theorem generalizes Fermat's: replace
	\(\mathbb Z_p^\times\) by \(\mathbb Z_n^*\). For example, one can show for all
	\(m \geq 3\) that \(\mathbb Z_{2^m}^*\) has exactly \(3\) elements of order \(2\), namely
	\([-1]\), \([1 + 2^{m-1}]\), and \([-(1 + 2^{m-1})]\). It follows that the product of
	all odd numbers \(r\) with \(1 \leq r < 2^m\) is congruent to \(1\) modulo \(2^m\), because
	\[
		(-1)(1 + 2^{m-1})(-1 - 2^{m-1}) = (1 + 2^{m-1})^2 = 1 + 2^m + 2^{2m-2}
		\equiv 1 \pmod{2^m}.
	\]
\end{remark}

\subsection{Quotient Groups}

The map \(\pi\colon \mathbb Z \to \mathbb Z_m\), \(\pi(a) = [a]_m\), is a surjective
homomorphism, so \(\mathbb Z_m = \im\pi\): every element of \(\mathbb Z_m\) has the form
\(\pi(a)\), and \(\pi(a) + \pi(b) = \pi(a+b)\). This description of \(\mathbb Z_m\) in
terms of \(\mathbb Z\) generalizes to arbitrary groups. Suppose \(f\colon G \to H\) is a
surjective homomorphism. Every element of \(H\) has the form \(f(a)\), and the operation
in \(H\) is given by \(f(a)f(b) = f(ab)\). Now \(K = \ker f\) is a normal subgroup of \(G\),
and we are going to reconstruct \(H = \im f\) (and a surjective homomorphism
\(\pi\colon G \to H\)) from \(G\) and \(K\) alone.

\begin{definition}[Product of subsets]
	Let \(\mathcal P(G)\) be the family of all nonempty subsets of a group \(G\). For
	\(X, Y \in \mathcal P(G)\), define
	\[
		XY = \{xy : x \in X \text{ and } y \in Y\}.
	\]
\end{definition}

This multiplication is associative: \(X(YZ)\) is the set of all \(x(yz)\) and \((XY)Z\) is
the set of all \((xy)z\) with \(x \in X\), \(y \in Y\), \(z \in Z\), and these are the same
elements. Some instances:
\begin{itemize}
	\item The product of a one-point subset \(\{a\}\) and a subgroup \(H\) is the coset
	      \(aH\).
	\item If \(H \leq G\), then \(HH = H\): \(HH \subseteq H\) because \(H\) is closed under
	      multiplication, and \(h = h1 \in HH\) for every \(h \in H\).
\end{itemize}
Two subsets may commute even though their elements do not: \(HH = HH\) for a nonabelian
subgroup \(H\). More interestingly, let \(G = S_3\) and \(K = \langle(1\ 2\ 3)\rangle\).
Then \((1\ 2)\) does not commute with \((1\ 2\ 3) \in K\), but \((1\ 2)K = K(1\ 2)\). In
fact, normal subgroups are characterized by this property.

\begin{lemma}\label{lem:normal-cosets}
	A subgroup \(K\) of a group \(G\) is normal if and only if \(bK = Kb\) for every
	\(b \in G\).
\end{lemma}

\begin{proof}
	Assume \(K \trianglelefteq G\), and let \(bk \in bK\). Then \(bkb^{-1} = k' \in K\), so
	\(bk = (bkb^{-1})b = k'b \in Kb\); hence \(bK \subseteq Kb\). For the reverse inclusion,
	let \(kb \in Kb\). Since \(K\) is normal, \(b^{-1}kb = b^{-1}k(b^{-1})^{-1} = k'' \in K\),
	so \(kb = b(b^{-1}kb) = bk'' \in bK\). Therefore \(bK = Kb\).

	Conversely (Rotman, Exercise 2.69), assume \(bK = Kb\) for all \(b\). If \(k \in K\) and
	\(b \in G\), then \(bk \in bK = Kb\), say \(bk = k'b\); hence
	\(bkb^{-1} = k' \in K\).
\end{proof}

Here is a fundamental construction of a new group from a given group.

\begin{theorem}\label{thm:quotient-group}
	Let \(G/K\) denote the family of all cosets of a subgroup \(K\) of \(G\). If \(K\) is a
	normal subgroup, then
	\[
		aK\,bK = abK
	\]
	for all \(a, b \in G\), and \(G/K\) is a group under this operation.
\end{theorem}

\begin{proof}
	The product \((aK)(bK)\) is a product of four elements of \(\mathcal P(G)\), so
	associativity in \(\mathcal P(G)\) gives
	\[
		(aK)(bK) = a(Kb)K = a(bK)K = ab(KK) = abK,
	\]
	using \(Kb = bK\) (\autoref{lem:normal-cosets}) and \(KK = K\). Thus the product of two
	cosets of \(K\) is again a coset of \(K\), and an operation on \(G/K\) has been defined.
	It is associative, because multiplication in \(\mathcal P(G)\) is. The identity is
	\(K = 1K\), for \((1K)(bK) = bK\), and the inverse of \(aK\) is \(a^{-1}K\), for
	\((a^{-1}K)(aK) = a^{-1}aK = K\). Therefore \(G/K\) is a group.
\end{proof}

\begin{definition}[Quotient group]
	The group \(G/K\) of \autoref{thm:quotient-group} is called the \vocab{quotient group}
	\(G\) mod \(K\). When \(G\) is finite, its order is
	\[
		\abs{G/K} = [G : K] = \abs{G}/\abs{K}
	\]
	(presumably the reason for the name).
\end{definition}

\begin{eg}[\(\mathbb Z/\langle m\rangle = \mathbb Z_m\)]\label{eg:Z-mod-m}
	Let \(\langle m\rangle = m\mathbb Z\) be the subgroup of all multiples of \(m\).
	Since \(\mathbb Z\) is abelian, \(\langle m\rangle\) is normal. The sets
	\(\mathbb Z/\langle m\rangle\) and \(\mathbb Z_m\) coincide, because they have the same
	elements: the coset \(a + \langle m\rangle = \{a + km : k \in \mathbb Z\}\) is the class
	\([a]_m\). The operations also coincide: addition in \(\mathbb Z/\langle m\rangle\) is
	\((a + \langle m\rangle) + (b + \langle m\rangle) = (a+b) + \langle m\rangle\), which is
	just \([a]_m + [b]_m = [a+b]_m\). Therefore \(\mathbb Z_m = \mathbb Z/\langle m\rangle\).
\end{eg}

Recall \autoref{lem:coset-basics}(i): if \(K \leq G\), then \(aK = bK\) if and only if
\(b^{-1}a \in K\); in particular \(aK = K\) if and only if \(a \in K\). We can now prove
the converse of \autoref{prop:kernel-image}(ii).

\begin{corollary}\label{cor:normal-is-kernel}
	Every normal subgroup \(K \trianglelefteq G\) is the kernel of some homomorphism.
\end{corollary}

\begin{proof}
	Define the \vocab{natural map} \(\pi\colon G \to G/K\) by \(\pi(a) = aK\). The formula
	\(aK\,bK = abK\) says \(\pi(a)\pi(b) = \pi(ab)\), so \(\pi\) is a (surjective)
	homomorphism. Since \(K\) is the identity of \(G/K\),
	\[
		\ker\pi = \{a \in G : aK = K\} = K. \qedhere
	\]
\end{proof}

The next theorem shows that every homomorphism gives rise to an isomorphism, and that
quotient groups are merely constructions of homomorphic images. E.~Noether (1882--1935)
emphasized the fundamental importance of this fact.

\begin{theorem}[First isomorphism theorem]\label{thm:first-iso}
	If \(f\colon G \to H\) is a homomorphism, then
	\[
		\ker f \trianglelefteq G \qquad\text{and}\qquad G/\ker f \cong \im f.
	\]
	In more detail, if \(\ker f = K\), then \(\varphi\colon G/K \to \im f \leq H\),
	\(\varphi(aK) = f(a)\), is an isomorphism.
\end{theorem}

\begin{proof}
	\(K = \ker f\) is normal by \autoref{prop:kernel-image}(ii).
	\begin{itemize}
		\item \emph{Well defined:} if \(aK = bK\), then \(a = bk\) for some \(k \in K\), and
		      \(f(a) = f(b)f(k) = f(b)\), because \(f(k) = 1\).
		\item \emph{Homomorphism:}
		      \(\varphi(aK\,bK) = \varphi(abK) = f(ab) = f(a)f(b) = \varphi(aK)\varphi(bK)\).
		\item \emph{Surjective onto \(\im f\):} clearly \(\im\varphi \subseteq \im f\); and if
		      \(y = f(a) \in \im f\), then \(y = \varphi(aK)\).
		\item \emph{Injective:} if \(\varphi(aK) = \varphi(bK)\), then \(f(a) = f(b)\), so
		      \(1 = f(b)^{-1}f(a) = f(b^{-1}a)\) and \(b^{-1}a \in K\); thus \(aK = bK\) by
		      \autoref{lem:coset-basics}(i). \qedhere
	\end{itemize}
\end{proof}

\begin{remark}
	The following commutative diagram describes the proof, where \(\pi\colon G \to G/K\) is
	the natural map; the theorem says \(f = \varphi\pi\) with \(\varphi\) an isomorphism
	onto \(\im f\).
	\[
		\begin{tikzcd}
			G \arrow[rr, "f"] \arrow[dr, "\pi"'] & & H \\
			& G/K \arrow[ur, "\varphi"'] &
		\end{tikzcd}
	\]
\end{remark}

For example, the identity \(1_G\colon G \to G\) is a surjective homomorphism with kernel
\(\{1\}\), so \(G/\{1\} \cong G\). Given any homomorphism \(f\colon G \to H\), one should
immediately ask for its kernel and image; the first isomorphism theorem then provides an
isomorphism \(G/\ker f \cong \im f\). Since there is no significant difference between
isomorphic groups, it also says that there is no significant difference between quotient
groups and homomorphic images.

\begin{eg}[Cyclic groups, again]\label{eg:cyclic-first-iso}
	We revisit \autoref{eg:cyclic-iso}. If \(G = \langle a\rangle\) is cyclic of order
	\(m\), define \(f\colon \mathbb Z \to G\) by \(f(n) = a^n\). Then \(f\) is a
	homomorphism (laws of exponents, \autoref{prop:exponents}), it is surjective because
	\(a\) generates \(G\), and \(\ker f = \{n \in \mathbb Z : a^n = 1\} = \langle m\rangle\)
	by \autoref{lem:order-divides}. The first isomorphism theorem gives
	\(\mathbb Z/\langle m\rangle \cong G\). So every cyclic group of order \(m\) is
	isomorphic to \(\mathbb Z/\langle m\rangle = \mathbb Z_m\), and hence any two cyclic
	groups of order \(m\) are isomorphic.
\end{eg}

\begin{eg}[\(\mathbb R/\mathbb Z \cong S^1\)]\label{eg:R-mod-Z}
	What is the quotient group \(\mathbb R/\mathbb Z\)? Define
	\(f\colon \mathbb R \to S^1\), where \(S^1\) is the circle group, by
	\(f(x) = e^{2\pi ix}\). Then \(f\) is a homomorphism, \(f(x+y) = f(x)f(y)\), by the
	addition formulas for sine and cosine, and \(f\) is surjective. Its kernel consists of
	all \(x \in \mathbb R\) with \(e^{2\pi ix} = \cos 2\pi x + i\sin 2\pi x = 1\); but
	\(\cos 2\pi x = 1\) forces \(x \in \mathbb Z\), and every integer lies in the kernel, so
	\(\ker f = \mathbb Z\). The first isomorphism theorem gives
	\[
		\mathbb R/\mathbb Z \cong S^1.
	\]
\end{eg}

\subsection{Products of Subgroups and the Isomorphism Theorems}

A natural question is whether \(HK\) is a subgroup when \(H\) and \(K\) are subgroups. In
general it is not: in \(G = S_3\), with \(H = \langle(1\ 2)\rangle\) and
\(K = \langle(1\ 3)\rangle\),
\[
	HK = \{(1), (1\ 2), (1\ 3), (1\ 2)(1\ 3)\} = \{(1), (1\ 2), (1\ 3), (1\ 3\ 2)\}
\]
has \(4\) elements, so it is not a subgroup lest we contradict Lagrange's theorem.

\begin{proposition}\label{prop:HK-subgroup}
	\begin{enumerate}[(i)]
		\item If \(H\) and \(K\) are subgroups of a group \(G\), and one of them is normal,
		      then \(HK\) is a subgroup of \(G\); moreover \(HK = KH\) in this case.
		\item If both \(H\) and \(K\) are normal subgroups, then \(HK\) is a normal subgroup.
	\end{enumerate}
\end{proposition}

\begin{proof}
	(i) Assume first that \(K \trianglelefteq G\). We claim \(HK = KH\). If \(hk \in HK\), then
	\(k' = hkh^{-1} \in K\), and \(hk = (hkh^{-1})h = k'h \in KH\); so \(HK \subseteq KH\).
	Similarly \(kh = h(h^{-1}kh) = hk'' \in HK\) with \(k'' = h^{-1}kh \in K\); so
	\(KH \subseteq HK\).

	Now \(HK\) is a subgroup: \(1 = 1 \cdot 1 \in HK\); if \(hk \in HK\), then
	\((hk)^{-1} = k^{-1}h^{-1} \in KH = HK\); and if \(hk, h_1k_1 \in HK\), then
	\(h_1^{-1}kh_1 = k' \in K\), so
	\[
		hk\,h_1k_1 = hh_1(h_1^{-1}kh_1)k_1 = (hh_1)(k'k_1) \in HK.
	\]
	If instead \(H \trianglelefteq G\), the same argument with the roles of \(H\) and \(K\)
	exchanged shows that \(KH = HK\) and that \(KH\) is a subgroup.

	(ii) If \(g \in G\), then \(ghkg^{-1} = (ghg^{-1})(gkg^{-1}) \in HK\).
\end{proof}

Here is a useful counting result.

\begin{proposition}[Product formula]\label{prop:product-formula}
	If \(H\) and \(K\) are subgroups of a finite group \(G\), then
	\[
		\abs{HK}\,\abs{H \cap K} = \abs{H}\,\abs{K},
	\]
	where \(HK = \{hk : h \in H,\ k \in K\}\) (not necessarily a subgroup).
\end{proposition}

\begin{proof}
	Define \(f\colon H \times K \to HK\) by \(f(h,k) = hk\); clearly \(f\) is surjective. Since
	\(H \times K\) is the disjoint union of the fibres
	\(f^{-1}(x) = \{(h,k) \in H \times K : hk = x\}\) for \(x \in HK\), it suffices to show
	\(\abs{f^{-1}(x)} = \abs{H \cap K}\) for every \(x \in HK\).

	We claim that if \(x = hk\), then
	\[
		f^{-1}(x) = \{(hd, d^{-1}k) : d \in H \cap K\}.
	\]
	Each \((hd, d^{-1}k)\) lies in \(f^{-1}(x)\), for \(hdd^{-1}k = hk = x\). Conversely, if
	\((h', k') \in f^{-1}(x)\), then \(h'k' = hk\), so \(h^{-1}h' = kk'^{-1}\) lies in
	\(H \cap K\); call it \(d\). Then \(h' = hd\) and \(k' = d^{-1}k\). Finally,
	\(d \mapsto (hd, d^{-1}k)\) is a bijection from \(H \cap K\) onto \(f^{-1}(x)\) (it is
	injective by cancellation), so \(\abs{f^{-1}(x)} = \abs{H \cap K}\).
\end{proof}

The next two results are variants of the first isomorphism theorem.

\begin{theorem}[Second isomorphism theorem]\label{thm:second-iso}
	If \(H\) and \(K\) are subgroups of a group \(G\) with \(H \trianglelefteq G\), then \(HK\)
	is a subgroup, \(H \cap K \trianglelefteq K\), and
	\[
		K/(H \cap K) \cong HK/H.
	\]
\end{theorem}

\begin{proof}
	\(HK\) is a subgroup by \autoref{prop:HK-subgroup}. Moreover \(H \trianglelefteq HK\): in
	general, if \(H \leq S \leq G\) and \(H\) is normal in \(G\), then \(H\) is normal in \(S\)
	(if \(ghg^{-1} \in H\) for all \(g \in G\), then in particular for all \(g \in S\)). So
	\(HK/H\) makes sense.

	Each coset \(xH \in HK/H\) has the form \(kH\) with \(k \in K\): write \(x = hk\); then
	\(hk = k(k^{-1}hk) = kh'\) with \(h' \in H\), so \(xH = kh'H = kH\). Hence the function
	\(f\colon K \to HK/H\), \(f(k) = kH\), is surjective. It is a homomorphism, being the
	restriction to \(K\) of the natural map \(HK \to HK/H\), and
	\[
		\ker f = \{k \in K : kH = H\} = K \cap H.
	\]
	Therefore \(H \cap K \trianglelefteq K\), and the first isomorphism theorem gives
	\(K/(H \cap K) \cong HK/H\).
\end{proof}

The second isomorphism theorem gives the product formula when one of the subgroups is
normal: \(\abs{K}/\abs{H \cap K} = \abs{HK}/\abs{H}\), that is,
\(\abs{HK}\,\abs{H \cap K} = \abs{H}\,\abs{K}\).

\begin{theorem}[Third isomorphism theorem]\label{thm:third-iso}
	If \(H\) and \(K\) are normal subgroups of a group \(G\) with \(K \leq H\), then
	\(H/K \trianglelefteq G/K\) and
	\[
		(G/K)/(H/K) \cong G/H.
	\]
\end{theorem}

\begin{proof}
	Define \(f\colon G/K \to G/H\) by \(f(aK) = aH\). It is well defined: if \(a'K = aK\),
	then \(a^{-1}a' \in K \leq H\), so \(a'H = aH\). It is a homomorphism,
	\(f(aK\,bK) = f(abK) = abH = aH\,bH\), and it is surjective. Moreover
	\[
		\ker f = \{aK : aH = H\} = \{aK : a \in H\} = H/K,
	\]
	so \(H/K\) is a normal subgroup of \(G/K\). Since \(f\) is surjective, the first
	isomorphism theorem gives \((G/K)/(H/K) \cong G/H\).
\end{proof}

The third isomorphism theorem is easy to remember: the \(K\)'s in the fraction
\((G/K)/(H/K)\) can be canceled. One better appreciates the first isomorphism theorem
after proving the third: \((G/K)/(H/K)\) consists of cosets (of \(H/K\)) whose
representatives are themselves cosets (of \(K\)), and a direct proof could be nasty.

The next result, which describes the subgroups of a quotient group \(G/K\), can be regarded
as a fourth isomorphism theorem. Recall that a function \(f\colon X \to Y\) sets up a
correspondence between subsets of \(X\) and subsets of \(Y\) via direct and inverse images;
we shall use the elementary facts
\[
	S \subseteq T \implies f(S) \subseteq f(T), \qquad
	S \subseteq f^{-1}(f(S)), \qquad
	f(f^{-1}(U)) = U \text{ if } f \text{ is surjective}.
\]
If \(K \trianglelefteq G\), let \(\operatorname{Sub}(G;K)\) be the family of all subgroups
\(S\) of \(G\) containing \(K\), and \(\operatorname{Sub}(G/K)\) the family of all subgroups
of \(G/K\).

\begin{proposition}[Correspondence theorem]\label{prop:correspondence}
	Let \(G\) be a group and \(K \trianglelefteq G\). Then \(S \mapsto S/K\) is a bijection
	\(\operatorname{Sub}(G;K) \to \operatorname{Sub}(G/K)\). Writing \(S^* = S/K\):
	\begin{enumerate}[(i)]
		\item \(T \leq S\) in \(\operatorname{Sub}(G;K)\) if and only if \(T^* \leq S^*\), in
		      which case \([S : T] = [S^* : T^*]\);
		\item \(T \trianglelefteq S\) in \(\operatorname{Sub}(G;K)\) if and only if
		      \(T^* \trianglelefteq S^*\), in which case \(S/T \cong S^*/T^*\).
	\end{enumerate}
\end{proposition}

\begin{proof}
	Let \(\pi\colon G \to G/K\) be the natural map, and note that
	\(S/K = \{sK : s \in S\} = \pi(S)\). This is a subgroup of \(G/K\): the image of a
	subgroup under a homomorphism is a subgroup (apply \autoref{prop:kernel-image}(i) to
	\(\pi\) restricted to \(S\)). Let \(\Phi(S) = S/K\).

	\emph{Key fact: if \(K \leq S \leq G\), then \(\pi^{-1}(\pi(S)) = S\).} Certainly
	\(S \subseteq \pi^{-1}(\pi(S))\). Conversely, if \(a \in \pi^{-1}(\pi(S))\), then
	\(\pi(a) = \pi(s)\) for some \(s \in S\); so \(s^{-1}a \in \ker\pi = K\), that is,
	\(a = sk\) with \(k \in K \leq S\), and \(a \in S\).

	\emph{\(\Phi\) is injective:} if \(\pi(S) = \pi(S')\) with \(S, S' \in
	\operatorname{Sub}(G;K)\), then \(S = \pi^{-1}\pi(S) = \pi^{-1}\pi(S') = S'\).

	\emph{\(\Phi\) is surjective:} if \(U \leq G/K\), then \(\pi^{-1}(U)\) is a subgroup of
	\(G\) containing \(\pi^{-1}(\{1\}) = K\) (\autoref{eg:kernels}(iv)), and
	\(\pi(\pi^{-1}(U)) = U\) because \(\pi\) is surjective.

	(i) If \(T \leq S\), then \(T^* = \pi(T) \subseteq \pi(S) = S^*\). Conversely, if
	\(T^* \leq S^*\), then \(T = \pi^{-1}(T^*) \subseteq \pi^{-1}(S^*) = S\) by the key fact.
	For the indices, we show that \(sT \mapsto \pi(s)T^*\) (\(s \in S\)) is a bijection from
	the cosets of \(T\) in \(S\) to the cosets of \(T^*\) in \(S^*\). For \(s, s' \in S\),
	\[
		sT = s'T \iff s'^{-1}s \in T = \pi^{-1}(T^*) \iff \pi(s')^{-1}\pi(s) \in T^*
		\iff \pi(s)T^* = \pi(s')T^*,
	\]
	which shows that the map is well defined and injective; it is surjective because every
	element of \(S^*\) is \(\pi(s)\) for some \(s \in S\). Hence \([S : T] = [S^* : T^*]\).
	(When \(G\) is finite this is simply
	\([S^*:T^*] = \frac{\abs{S}/\abs{K}}{\abs{T}/\abs{K}} = \abs{S}/\abs{T} = [S:T]\).)

	(ii) If \(T \trianglelefteq S\), then \(K \leq T\) are normal subgroups of the group
	\(S\), and the third isomorphism theorem (applied to \(S\)) gives
	\(T/K \trianglelefteq S/K\) and \((S/K)/(T/K) \cong S/T\), that is,
	\(S^*/T^* \cong S/T\). Conversely, assume \(T^* \trianglelefteq S^*\), and let
	\(t \in T\), \(s \in S\). Then
	\[
		\pi(sts^{-1}) = \pi(s)\pi(t)\pi(s)^{-1} \in \pi(s)T^*\pi(s)^{-1} = T^*,
	\]
	so \(sts^{-1} \in \pi^{-1}(T^*) = T\). Hence \(T \trianglelefteq S\).
\end{proof}

When dealing with quotient groups, one usually says, without explicitly mentioning the
correspondence theorem, that every subgroup of \(G/K\) has the form \(S/K\) for a unique
subgroup \(S\) with \(K \leq S \leq G\).

\begin{eg}[The second center]
	If \(G\) is a finite \(p\)-group and \(G \neq \{1\}\), then \(Z(G) \neq \{1\}\) by
	\autoref{thm:p-group-center}. Hence, if \(G\) is not abelian, then \(G/Z(G)\) is a
	nontrivial \(p\)-group and \(Z(G/Z(G)) \neq \{1\}\). An important role in the study of
	finite \(p\)-groups is played by \(Z_2(G)\), defined as the inverse image of
	\(Z(G/Z(G))\) under the natural map \(G \to G/Z(G)\). Since the center
	\(Z(G/Z(G))\) is normal in \(G/Z(G)\), the correspondence theorem gives
	\(Z_2(G) \trianglelefteq G\).
\end{eg}

\begin{proposition}\label{prop:abelian-subgroups}
	If \(G\) is a finite abelian group, then \(G\) has a subgroup of order \(d\) for every
	divisor \(d\) of \(\abs{G}\). In particular, if \(p\) is a prime divisor of \(\abs{G}\),
	then \(G\) contains an element of order \(p\). (Compare \autoref{prop:A4-no-6}.)
\end{proposition}

\begin{proof}
	First we prove, by induction on \(n = \abs{G}\), that for every prime \(p \mid n\) there
	is an element of order \(p\) in \(G\). The base step \(n = 1\) is vacuous. For the
	inductive step, \(n > 1\), so we may choose \(a \in G\) of order \(k > 1\).

	If \(p \mid k\), say \(k = p\ell\), then \(a^\ell\) has order \(p\): indeed
	\((a^\ell)^p = a^k = 1\), while \((a^\ell)^j = a^{\ell j} \neq 1\) for \(0 < j < p\),
	since \(0 < \ell j < k\).

	If \(p \nmid k\), let \(H = \langle a\rangle\). Since \(G\) is abelian,
	\(H \trianglelefteq G\), and \(\abs{G/H} = n/k\) is divisible by \(p\) (Euclid's lemma).
	By induction, \(G/H\) has an element \(bH\) of order \(p\). If \(b\) has order \(m\),
	then \((bH)^m = b^mH = H\), so \(p \mid m\) by \autoref{lem:order-divides}. We are back
	in the first case: \(b^{m/p}\) has order \(p\).

	Now we prove the general statement by induction on \(d \geq 1\), for all finite abelian
	groups at once. The base step \(d = 1\) is obvious. If \(d > 1\), let \(p\) be a prime
	divisor of \(d\). By the first part, \(G\) has a subgroup \(H\) of order \(p\); it is
	normal because \(G\) is abelian, and \(\abs{G/H} = \abs{G}/p\) is divisible by \(d/p\).
	By induction, \(G/H\) has a subgroup \(S^*\) of order \(d/p\). By the correspondence
	theorem, \(S^* = S/H\) for a subgroup \(S\) with \(H \leq S \leq G\), and
	\(\abs{S} = \abs{S^*}\,\abs{H} = (d/p) \cdot p = d\).
\end{proof}

Cauchy's theorem (\autoref{thm:cauchy}) says that the second statement holds for every
finite group, abelian or not.

\subsection{Direct Products}

Here is another construction of a new group from two given groups.

\begin{definition}[Direct product]\label{def:direct-product}
	If \(H\) and \(K\) are groups, their \vocab{direct product} \(H \times K\) is the set of
	all ordered pairs \((h, k)\) with \(h \in H\) and \(k \in K\), with the operation
	\[
		(h, k)(h', k') = (hh', kk').
	\]
\end{definition}

It is routine to check that \(H \times K\) is a group: associativity holds coordinatewise,
the identity is \((1, 1)\), and \((h, k)^{-1} = (h^{-1}, k^{-1})\). Note that \(H \times K\)
is abelian if and only if both \(H\) and \(K\) are abelian.

\begin{eg}[\(\mathbf V \cong \mathbb Z_2 \times \mathbb Z_2\)]\label{eg:V-Z2Z2}
	The function \(f\colon \mathbf V \to \mathbb Z_2 \times \mathbb Z_2\) defined by
	\begin{align*}
		f\colon (1)          & \mapsto ([0], [0]), & (1\ 2)(3\ 4) & \mapsto ([1], [0]), \\
		(1\ 3)(2\ 4)         & \mapsto ([0], [1]), & (1\ 4)(2\ 3) & \mapsto ([1], [1])
	\end{align*}
	is an isomorphism: in both groups, every nonidentity element has order \(2\) and the
	product of any two distinct nonidentity elements is the third one.
\end{eg}

We now apply the first isomorphism theorem to direct products.

\begin{proposition}\label{prop:product-quotient}
	Let \(G\) and \(G'\) be groups, and let \(K \trianglelefteq G\) and
	\(K' \trianglelefteq G'\). Then \(K \times K'\) is a normal subgroup of \(G \times G'\),
	and
	\[
		(G \times G')/(K \times K') \cong (G/K) \times (G'/K').
	\]
\end{proposition}

\begin{proof}
	Let \(\pi\colon G \to G/K\) and \(\pi'\colon G' \to G'/K'\) be the natural maps, and
	define \(f\colon G \times G' \to (G/K) \times (G'/K')\) by
	\[
		f(g, g') = (\pi(g), \pi'(g')) = (gK, g'K').
	\]
	Since \(\pi\) and \(\pi'\) are surjective homomorphisms and the operations are
	coordinatewise, \(f\) is a surjective homomorphism. Its kernel is
	\(\{(g, g') : gK = K,\ g'K' = K'\} = K \times K'\). The first isomorphism theorem now
	gives both assertions.
\end{proof}

Here is a characterization of direct products.

\begin{proposition}\label{prop:internal-direct}
	If a group \(G\) contains normal subgroups \(H\) and \(K\) with \(H \cap K = \{1\}\) and
	\(HK = G\), then \(G \cong H \times K\).
\end{proposition}

\begin{proof}
	First, each \(g \in G\) has a \emph{unique} factorization \(g = hk\) with \(h \in H\) and
	\(k \in K\): if \(hk = h'k'\), then \(h'^{-1}h = k'k^{-1} \in H \cap K = \{1\}\), so
	\(h' = h\) and \(k' = k\). We may therefore define \(\varphi\colon G \to H \times K\) by
	\(\varphi(g) = (h, k)\), where \(g = hk\).

	Next, elements of \(H\) commute with elements of \(K\). Let \(h \in H\) and \(k \in K\).
	Since \(K\) is normal, \((hkh^{-1})k^{-1} \in K\); since \(H\) is normal,
	\(h(kh^{-1}k^{-1}) \in H\). But \(H \cap K = \{1\}\), so \(hkh^{-1}k^{-1} = 1\), and
	\(hk = kh\).

	Now \(\varphi\) is a homomorphism: if \(g = hk\) and \(g' = h'k'\), then
	\(gg' = hkh'k' = (hh')(kk')\), so
	\[
		\varphi(gg') = (hh', kk') = (h, k)(h', k') = \varphi(g)\varphi(g').
	\]
	It is surjective, since \(\varphi(hk) = (h, k)\), and injective, since
	\(\varphi(g) = (1, 1)\) forces \(g = 1\). Therefore \(\varphi\) is an isomorphism.
\end{proof}

All the hypotheses are needed. For example, let \(G = S_3\), \(H = \langle(1\ 2\ 3)\rangle\),
and \(K = \langle(1\ 2)\rangle\). Then \(S_3 = HK\) (by the product formula,
\(\abs{HK} = 3 \cdot 2 / 1 = 6\)), \(H \cap K = \{1\}\), and \(H \trianglelefteq S_3\), but
\(K\) is not normal. And \(S_3 \not\cong H \times K\), for \(S_3\) is not abelian, while a
direct product of abelian groups is abelian.

\begin{theorem}\label{thm:Zmn}
	If \(m\) and \(n\) are relatively prime, then
	\[
		\mathbb Z_{mn} \cong \mathbb Z_m \times \mathbb Z_n.
	\]
\end{theorem}

\begin{proof}
	Define \(f\colon \mathbb Z \to \mathbb Z_m \times \mathbb Z_n\) by
	\(f(a) = ([a]_m, [a]_n)\); it is a homomorphism because each coordinate is. We claim
	\(\ker f = \langle mn\rangle\). Clearly \(\langle mn\rangle \leq \ker f\). Conversely, if
	\(a \in \ker f\), then \([a]_m = [0]_m\) and \([a]_n = [0]_n\), so \(m \mid a\) and
	\(n \mid a\); since \(\gcd(m, n) = 1\), \autoref{lem:coprime-divisibility}(ii) gives
	\(mn \mid a\), so \(a \in \langle mn\rangle\).

	By the first isomorphism theorem and \autoref{eg:Z-mod-m},
	\(\mathbb Z_{mn} = \mathbb Z/\langle mn\rangle \cong \im f\). Thus \(\im f\) is a
	subgroup of \(\mathbb Z_m \times \mathbb Z_n\) with \(mn\) elements; since
	\(\mathbb Z_m \times \mathbb Z_n\) also has \(mn\) elements, \(\im f\) is all of it,
	and \(\mathbb Z_{mn} \cong \mathbb Z_m \times \mathbb Z_n\).
\end{proof}

\begin{remark}
	Surjectivity of \(f\) says: for any \(a, b \in \mathbb Z\) there is \(x \in \mathbb Z\)
	with \(x \equiv a \pmod m\) and \(x \equiv b \pmod n\). This is the \vocab{Chinese
		remainder theorem}; Rotman uses it to prove surjectivity, while the counting argument
	above proves it.
\end{remark}

For example, \(\mathbb Z_6 \cong \mathbb Z_2 \times \mathbb Z_3\). There is no such
isomorphism if \(m\) and \(n\) are not relatively prime: \autoref{eg:V-Z2Z2} shows
\(\mathbf V \cong \mathbb Z_2 \times \mathbb Z_2\), which is not isomorphic to
\(\mathbb Z_4\), because \(\mathbf V\) has no element of order \(4\).

Since \(\langle a\rangle \cong \mathbb Z_n\) exactly when \(a\) has order \(n\)
(\autoref{prop:order-equals-size} and \autoref{prop:Zm-cyclic}), \autoref{thm:Zmn} can be read as saying that if \(a\) and \(b\)
are commuting elements of relatively prime orders \(m\) and \(n\), then \(ab\) has order
\(mn\). Here is a direct proof (this is the statement we already used as
\autoref{lem:coprime-orders}).

\begin{proposition}\label{prop:coprime-orders}
	Let \(G\) be a group, and let \(a, b \in G\) be commuting elements of orders \(m\) and
	\(n\). If \(\gcd(m, n) = 1\), then \(ab\) has order \(mn\).
\end{proposition}

\begin{proof}
	Since \(a\) and \(b\) commute, \((ab)^r = a^rb^r\) for all \(r\), so
	\((ab)^{mn} = a^{mn}b^{mn} = 1\). It suffices to prove that \((ab)^k = 1\) implies
	\(mn \mid k\). If \(1 = (ab)^k = a^kb^k\), then \(a^k = b^{-k}\). Since \(a\) has order
	\(m\), \(1 = a^{mk} = b^{-mk}\); since \(b\) has order \(n\), \autoref{lem:order-divides}
	gives \(n \mid mk\), and so \(n \mid k\) by Euclid's lemma. Similarly \(m \mid k\).
	Finally \(mn \mid k\) by \autoref{lem:coprime-divisibility}(ii).
\end{proof}

Here is a number-theoretic application of direct products.

\begin{corollary}\label{cor:phi-multiplicative}
	If \(\gcd(m, n) = 1\), then \(\phi(mn) = \phi(m)\phi(n)\).
\end{corollary}

\begin{proof}
	The isomorphism of \autoref{thm:Zmn} is
	\(f\colon \mathbb Z_{mn} \to \mathbb Z_m \times \mathbb Z_n\),
	\(f([a]_{mn}) = ([a]_m, [a]_n)\) (this is the \(\varphi\) of the first isomorphism
	theorem). It also preserves multiplication:
	\[
		f([a]_{mn}[b]_{mn}) = f([ab]_{mn}) = ([a]_m[b]_m, [a]_n[b]_n) = f([a]_{mn})f([b]_{mn}),
	\]
	where pairs are multiplied coordinatewise, and \(f([1]_{mn}) = ([1]_m, [1]_n)\).

	We claim \(f(\mathbb Z_{mn}^*) = \mathbb Z_m^* \times \mathbb Z_n^*\). If
	\([a]_{mn}[b]_{mn} = [1]_{mn}\), then applying \(f\) gives \([a]_m[b]_m = [1]_m\) and
	\([a]_n[b]_n = [1]_n\), so \(f([a]_{mn}) \in \mathbb Z_m^* \times \mathbb Z_n^*\).
	Conversely, suppose \(f([c]_{mn}) = ([c]_m, [c]_n)\) with \([c]_m[d]_m = [1]_m\) and
	\([c]_n[e]_n = [1]_n\). Since \(f\) is surjective, there is \(b\) with
	\(([b]_m, [b]_n) = ([d]_m, [e]_n)\); then
	\(f([c]_{mn}[b]_{mn}) = ([1]_m, [1]_n) = f([1]_{mn})\), and injectivity gives
	\([c]_{mn}[b]_{mn} = [1]_{mn}\), so \([c]_{mn} \in \mathbb Z_{mn}^*\).

	By \autoref{lem:units-order},
	\(\phi(mn) = \abs{\mathbb Z_{mn}^*} = \abs{\mathbb Z_m^* \times \mathbb Z_n^*}
	= \phi(m)\phi(n)\).
\end{proof}

\begin{definition}[Direct product of \(n\) groups]
	If \(H_1, \dots, H_n\) are groups, their \vocab{direct product}
	\(H_1 \times \cdots \times H_n\) is the set of all \(n\)-tuples \((h_1, \dots, h_n)\)
	with \(h_i \in H_i\), with coordinatewise multiplication
	\[
		(h_1, \dots, h_n)(h_1', \dots, h_n') = (h_1h_1', \dots, h_nh_n').
	\]
\end{definition}

The \emph{basis theorem} (Rotman, Theorem 6.11) says that every finite abelian group is a
direct product of cyclic groups. We end with a variant of \autoref{prop:cyclic-criterion}.

\begin{proposition}\label{prop:unique-p-subgroup-cyclic}
	If \(G\) is a finite abelian group having a unique subgroup of order \(p\) for every prime
	divisor \(p\) of \(\abs{G}\), then \(G\) is cyclic.
\end{proposition}

\begin{proof}
	Choose \(a \in G\) of largest order, say \(n\). We first show:
	\emph{\(\langle a\rangle\) contains every element \(x \in G\) with \(x^p = 1\) for some
		prime \(p\).} If \(x \neq 1\), then \(x\) has order \(p\), so \(p \mid \abs{G}\) and
	\(C = \langle x\rangle\) is the unique subgroup of order \(p\). If \(\gcd(p, n) = 1\),
	then \(xa\) has order \(pn > n\) by \autoref{prop:coprime-orders}, contradicting the
	choice of \(a\). So \(p \mid n\), say \(n = pq\); then \(a^q\) has order \(p\), so
	\(\langle a^q\rangle = C\) by uniqueness, and \(x \in C \leq \langle a\rangle\).

	If \(\langle a\rangle = G\), we are done; so assume there is \(b \in G\) with
	\(b \notin \langle a\rangle\). Since \(b^{\abs{G}} = 1 \in \langle a\rangle\), there is a
	smallest positive integer \(k\) with \(b^k \in \langle a\rangle\), say
	\[
		b^k = a^r.
	\]
	Now \(k\) is the order of \(b\langle a\rangle\) in \(G/\langle a\rangle\), so
	\(k \mid \abs{G}\); and \(k \neq 1\). Write \(k = p\ell\) with \(p\) prime; then
	\(p \mid \abs{G}\). By \autoref{prop:abelian-subgroups}, \(G\) has an element of order
	\(p\), and the argument of the first paragraph shows that \(p \mid n\). There are two
	possibilities.
	\begin{itemize}
		\item If \(p \mid r\), say \(r = pu\), then \(b^{p\ell} = b^k = a^{pu}\), so
		      \((b^\ell a^{-u})^p = 1\) (as \(G\) is abelian). By the first paragraph,
		      \(b^\ell a^{-u} \in \langle a\rangle\), so \(b^\ell \in \langle a\rangle\) with
		      \(0 < \ell < k\), contradicting the minimality of \(k\).
		\item If \(p \nmid r\), then \(\gcd(p, r) = 1\), so \(1 = sp + tr\) for some
		      integers \(s, t\), and
		      \[
			      a = a^{sp + tr} = a^{sp}(a^r)^t = a^{sp}b^{p\ell t} = (a^sb^{\ell t})^p.
		      \]
		      So \(a = x^p\) with \(x = a^sb^{\ell t}\). Let \(x\) have order \(N\). Then
		      \(a^N = (x^N)^p = 1\), so \(n \mid N\). If \(N = n\), then, since \(p \mid n\),
		      the element \(x^p\) would have order \(n/p\) (as \((x^p)^{n/p} = 1\) and
		      \((x^p)^j = x^{pj} \neq 1\) for \(0 < j < n/p\)), contradicting
		      \(x^p = a\) of order \(n\). Hence \(N > n\), contradicting the choice of \(a\).
	\end{itemize}
	Both cases are impossible, so \(G = \langle a\rangle\).
\end{proof}

The proposition is false for nonabelian groups: the quaternions \(\mathbf Q\) form a
noncyclic group of order \(8\) having a unique subgroup of order \(2\)
(\autoref{eg:Q-subgroups-normal}).
